\chapter*{Presentación}
\addcontentsline{toc}{chapter}{Presentación}
\markboth{Presentación}{}
Este tercer volumen estudia cómo las propiedades geométricas y topológicas de un espacio controlan sus funcionales y operadores. El recorrido comienza con extensión y separación, pasa por el espectro y las topologías débiles, y culmina en álgebras de Banach, álgebras C* y álgebras de von Neumann. Cada etapa prepara herramientas que la siguiente utiliza: la dualidad permite separar; la compacidad produce espectros discretos; las representaciones convierten elementos abstractos en operadores.

La fuente principal son las lecturas 0--36 de \emph{Análisis Matemático 3. Licenciatura en Física y Matemáticas}, de Miguel Ángel Valencia Bucio \cite{valencia3}. La edición, organización temática, explicaciones ampliadas, soluciones y rectificaciones corresponden a Eder Said Martinez Ramirez. Cada bloque identifica la lectura de origen; la tabla final relaciona cada ejercicio con su número original y enlaza los archivos consultados. Los ejercicios complementarios se distinguen de los originales.

El libro está dirigido a estudiantes que conocen álgebra lineal, cálculo, espacios métricos y normados y los fundamentos de la integral de Lebesgue. Los volúmenes I y II de esta colección proporcionan esos antecedentes. La lectura preliminar repasa sus conexiones y corrige algunas formulaciones que, tomadas literalmente, serían falsas. Los argumentos complejos utilizan el teorema de Cauchy, series de potencias, el principio de identidad y Liouville; se explica cómo pasan de funciones escalares a expresiones con valores en operadores.

\section*{Apuntes y solucionario}
Los ejercicios se presentan en estos apuntes sin soluciones. Los mismos
enunciados, con la misma numeración por capítulo y sus soluciones, se reúnen
en el documento separado \emph{Solucionario} de este volumen. Las menciones
a las soluciones y a los ejercicios resueltos de esta presentación se refieren
a ese documento complementario.

\section*{Cómo leer y verificar una demostración}
Antes de calcular, identifique el espacio donde vive cada objeto y la topología utilizada. Una sucesión puede converger débilmente y no en norma; un adjunto entre duales complejos es lineal en el operador, mientras que el adjunto de Hilbert es conjugado lineal. Un operador puede tener inversa algebraica sin una inversa continua si falta completitud. Estas distinciones intervienen en las pruebas, no son convenciones decorativas.

Una demostración se considera completa cuando explica las elecciones, comprueba las hipótesis y justifica el paso al límite. Si el enunciado original es falso, la solución ofrece un contraejemplo verificable y, cuando corresponde, una versión corregida con demostración. No se presenta una afirmación distinta como si fuera la original. Se registran por separado las correcciones matemáticas y las pérdidas de signos debidas a la extracción del PDF.

\section*{Convenciones y fuentes complementarias}
El producto interno es lineal en la primera variable. Escribimos $\Bcal(X,Y)$ para los operadores lineales acotados y $X^*$ para el dual continuo. En los capítulos de álgebras, $\Lcal(H)$ designa también $\Bcal(H)$. Usamos $T'$ para el adjunto de Banach y $T^*$ para el adjunto de Hilbert o la involución de un álgebra C*. El espectro se estudia en espacios complejos no nulos y en álgebras unitarias no nulas, salvo que se indique expresamente otra convención. En un álgebra sin unidad se utiliza la unitalización que se construye en el texto.

Conway \cite{conway1990}, Kreyszig \cite{kreyszig1978} y Rudin \cite{rudin1991} aportan el contexto del análisis funcional. Murphy \cite{murphy1990} y Zhu \cite{zhu1993} desarrollan las álgebras de operadores. Las notas de Egor Maximenko se citan localmente en las ampliaciones sobre Schur, desplazamientos, descomposición espectral, Gelfand y cálculo continuo. La presentación de operadores de multiplicación pertenece a Gerardo Ramos Vázquez \cite{ramos2018}; que esté alojada en el sitio de otro profesor no modifica su autoría.

\section*{Créditos del diseño}
El diseño toma como referencia visual los \emph{Apuntes de Homotopía Estable} del repositorio público de MrPaborru (\url{https://github.com/MrPaborru/Apuntes-Homotop-a-Estable}), combinados con una composición tradicional de libros de matemáticas. El emblema de la ESFM procede de su sitio oficial. Su uso identifica la adscripción académica del autor; esta colección es una edición personal, no una publicación oficial ni una obra avalada por el IPN. Se conserva una paleta inspirada en los colores de la ESFM y una tipografía Latin Modern. Los detalles de procedencia gráfica se documentan también en las fuentes LaTeX.
